\subsection{An exact return-prime cover retaining the entire old pair}

\begin{statement}{D.46}{Proved}{A strict mixed return cover with its complete generating peripheral fibre}
Let \(\ell,p\geq5\) be distinct primes, put
\(a=\operatorname{ord}_p(\ell)\), and choose
\(k=\mathbf F_{\ell^a}\) with an element \(\zeta\) of order \(p\).
Indices below lie in \(\mathbf Z/\ell\).
Let \(c\) cyclically permute coordinates, put
\[
 N=(\mathbf F_p)^\ell,\quad
 P=N\rtimes\langle c\rangle,\quad
 V=k^\ell,\quad H=V\rtimes P,
 \qquad
 n(v_i)_i=(\zeta^{n_i}v_i)_i.
 \tag{D.46.1}
\]
These are specified actions, and the action of \(c\) on both
\(N\) and \(V\) is the same cyclic shift.
Retain the actual source \(\widehat B,F,s,r,\tau,x,y,z\) of D.44.
There is an epimorphism
\[
 q_H:F\twoheadrightarrow H,
 \qquad x\longmapsto X=(e_0,c),\quad
 y\longmapsto Y=(0,n_0),\quad z\longmapsto Z=(XY)^{-1},
 \tag{D.46.2}
\]
where \(n_0\) is the first coordinate vector of \(N\), and
\(e_0\) is the first coordinate vector of \(V\).
Let \(E\to D\) be the actual six-coset induced images of
\(\widehat B\) obtained from \(q_H\) and its projection
\(q_P:F\to P\), respectively, as specified in
\textup{(D.46.5)} below. Then
\[
 1\longrightarrow W=V^6\longrightarrow E\longrightarrow D
 \longrightarrow1,
 \qquad \dim_{\mathbf F_\ell}W=d=6\ell a.
 \tag{D.46.3}
\]
The \(D\)-module \(W\) is irreducible and nontrivial.

Put \(K=O^p(O^\ell(F))\) and
\(\Gamma=\widehat B/K\).
Every D.44 normalized automorphism of this entire \(\Gamma\)
induces an automorphism of the specified \(D\), and has a lift
to the actual joint source
\[
 J=\Gamma\mathbin{\times_D}E
   =\widehat B/(K\cap\ker q_E).
 \tag{D.46.4}
\]
The lift fixes \(\tau\), retains the complete old normalized
conjugator in \(\Gamma\), and has its new conjugator in the
image of \(F'\). Thus it retains the entire R8 pro-\(\ell\)
pair and the chosen whole relative pro-\(p\) lift of D.44.
The maximal pro-\(\ell\) quotient of the pure subgroup of
\(J\) is exactly \(P_\ell\), with the original action.

For the specified old automorphism, the complete literal
peripheral fibre before fixing \(\tau\) has \(\ell^d\)
elements. Every one generates the actual joint source and
factors through an automorphism of that source.
This fibre is one free \(W\)-conjugacy orbit.
The same set is obtained with any prescribed three base
conjugators, including the projections of prescribed old
conjugators in \(\Gamma\).
Fixing \(\tau\) gives precisely
\(\ell^{\ell a}\) literal normalized maps, a free transitive
orbit of \(W^\tau\). These counts concern maps and their
conjugator images; the complete profinite word fibres are
the cosets described in \textup{(D.46.17)}.

The source in \textup{(D.46.4)} is an actual quotient of
the return cover
\(\widehat B/O^\ell(O^p(O^\ell(F)))\), with a
nontrivial finite kernel over \(\Gamma\).
In contrast, adjacent repetition of one residual has the
exact identity \(O^p(O^p L)=O^p L\) for every profinite
group \(L\).
\end{statement}

\begin{proof}
We give the source construction and all descent maps explicitly.
Choose the six representatives \(t_i\) of \(B/F\) used in
D.44, with \(t_0=1\). For \(b\in\widehat B\), let
\(j=i\cdot\bar b\) in the regular right action of
\(\mathfrak S_3\), and put
\[
 q_E(b)=\left((i\mapsto j),
           \bigl(q_H(t_i b t_j^{-1})\bigr)_i\right),
 \qquad
 q_D(b)=\left((i\mapsto j),
           \bigl(q_P(t_i b t_j^{-1})\bigr)_i\right).
 \tag{D.46.5}
\]
The formula is first evaluated on the two literal generators
\(s,r\), and gives their continuous maps from the profinite
free product. In fact
\(t_i b b' t_k^{-1}=(t_i b t_j^{-1})(t_j b't_k^{-1})\);
this proves the multiplication rule, the two source relations,
and the displayed projections. The kernels are the corresponding
normal cores of \(\ker q_H\) and \(\ker q_P\) in
\(\widehat B\). Both images retain the specified
\(\mathfrak S_3\) quotient. No additional quotient kernel
is being presumed invariant under a changed source map.

\paragraph{The two-generator group and its residuals.}
The images of \(X,Y\) generate \(P\): conjugating \(n_0\)
by powers of \(c\) gives every coordinate generator of \(N\).
Moreover \(X^\ell=\sum_i e_i\ne0\) in \(V\).
Acting on this vector by \(n_0-1\) gives
\((\zeta-1)e_0\); powers of \(n_0\), their
\(\mathbf F_\ell\)-linear combinations, and the cyclic
shift generate all of \(V\), because
\(\mathbf F_\ell[\zeta]=k\).
This proves the epimorphism \textup{(D.46.2)}.
The same calculation proves that \(V\) is simple as an
\(\mathbf F_\ell[P]\)-module. Indeed the distinct
coordinate characters of the prime-to-\(\ell\) subgroup
\(N\) give its \(\ell\) simple coordinate summands
over \(\mathbf F_\ell\); \(c\) permutes them transitively.

The following are equalities with their quotient maps:
\[
 O^\ell(H)=V\rtimes N,
 \qquad O^p(V\rtimes N)=V,
 \qquad O^\ell(V)=1,
 \qquad O^\ell(P)=N,\quad O^p(N)=1.
 \tag{D.46.6}
\]
For the first equality, every finite \(\ell\)-quotient kills
\(N\), and therefore kills \([V,N]=V\); the remaining
quotient is the displayed \(C_\ell\). For the second,
every finite \(p\)-quotient kills \(V\), and the quotient
by \(V\) is the displayed \(p\)-group \(N\).
The last two identities follow by the same source projections.

The action of \(X\) on \(V\) is the cyclic permutation,
the action of \(Y\) changes only the first coordinate by
\(\zeta\), and the action of \(XY\) has its \(\ell\)-th
power equal to \(\zeta I\). Consequently
\[
 \begin{array}{c|ccc}
       &X&Y&Z\\ \hline
 \text{order in }H&\ell^2&p&\ell p\\
 \text{order in }P&\ell&p&\ell p\\
 \dim_{\mathbf F_\ell}V^{(\cdot)}&a&(\ell-1)a&0.
 \end{array}
 \tag{D.46.7}
\]
For \(XY\), the operator \(XY-1\) is invertible since
\((XY)^\ell=\zeta I\); hence the translation part of
\((XY)^{\ell p}\) vanishes, giving its asserted order.

\paragraph{The entire induced kernel.}
The kernel of \(E\to D\) is a submodule of \(V^6\).
The restrictions of its six ambient summands to the image of
\(F\) are simple. Their fixed-space profiles on the same
\(x,y,z\) are the six permutations of
\((a,(\ell-1)a,0)\), by the actual peripheral
\(\mathfrak S_3\) permutations in D.44.
They are therefore pairwise nonisomorphic.
For every summand one of
\(x^{\ell p},y^{\ell p},z^{\ell p}\) has nonzero
projection in that summand and is killed in \(D\):
use the \(X\) coordinate in \textup{(D.46.7)}.
Its normal closure surjects onto that simple summand.
A submodule of a direct sum of pairwise nonisomorphic simple
modules which surjects onto each summand is the whole direct
sum. This follows, for example, by splitting a submodule of
the semisimple module and using the absence of homomorphisms
between distinct simple summands; it does not require the
whole ambient group algebra to be semisimple.
Thus the kernel is precisely \(W=V^6\).
The \(\mathfrak S_3\) quotient permutes these six simple
summands transitively, so \(W\) is simple for \(D\).
In particular
\[
 W^D=0,
 \qquad W_D=0,
 \qquad [E_F,W]=W\subset E_F',
 \quad E_F=\operatorname{im}(F\to E).
 \tag{D.46.8}
\]

The action of \(s\) on the six summands has three two-cycles,
and the action of \(r\) has two three-cycles. Their powers
are literally the identity, so their invariant dimensions
are \(d/2,d/3\). The action of \(\tau\) has three
two-cycles and \(\tau^2=x^{-1}\). On a two-cycle its
fixed vectors are exactly the fixed vectors of \(x\) on
one chosen summand. Formula \textup{(D.46.7)} gives
\[
 \dim W^s=3\ell a,\qquad
 \dim W^r=2\ell a,\qquad
 \dim W^\tau=\ell a.
 \tag{D.46.9}
\]
Also the orders of \(\tau\) in \(D,E\) are respectively
\(2\ell p,2\ell^2p\), by \(\tau^2=x^{-1}\) and
\textup{(D.46.7)}.

\paragraph{All centralizer projections and all conjugator choices.}
For \(s,r\), the maps from the upper centralizers to the
lower centralizers are onto. A lift of a lower centralizer
element conjugates the chosen upper element to another
lift of the same cyclic subgroup; averaging over its order
\(2\) or \(3\), which is invertible in
\(\mathbf F_\ell\), supplies the required kernel
conjugator. Equivalently, this is the explicit coprime
cyclic-complement argument of D.44.

For \(\tau\) we verify the extra assertion directly.
The local centralizer projections are
\[
 \begin{split}
 C_H(X)&\longrightarrow\langle c\rangle
                    \subset C_P(c)=\langle c\rangle\times Z(P),\\
 C_H(Y)&\twoheadrightarrow C_P(n_0)=N,\\
 C_H(Z)&\twoheadrightarrow C_P(q_P(z)).
 \end{split}
 \tag{D.46.10}
\]
The image in the first line is exactly \(\langle c\rangle\).
Indeed an element of \(C_P(c)\) has its \(N\)-coordinate
constant, say \(b\); summing the coordinates of the
translation equation for commuting with \(X\) gives
\(1=\zeta^b\), hence \(b=0\). Powers of \(X\)
give every power of \(c\). The second line has the
literal commuting lifts \((0,n)\). For the third,
\(1-Z\) is invertible on \(V\); the translation
equation therefore has one solution for each lower
centralizer element.

The pure image in \(D\) retains the actual quotient
\(F^{\rm ab}/p\simeq\mathbf F_p^2\).
The second abelian coordinate is the sum of the
\(N\)-coordinates in the identity component of
\textup{(D.46.5)}, and conjugating by \(s\) recovers
the first. Conjugation by \(\tau\) sends these coordinates
\((u,v)\) to \((u-v,-v)\), since it fixes \(x\)
and interchanges the other two peripheral classes.
An element of the pure image centralizing \(\tau\)
therefore has \(v=0\). In the identity component its
projection to \(C_P(c)\) has constant \(N\)-coordinate
\(b\) with \(\ell b=0\) in \(\mathbf F_p\), so
\(b=0\). Its projection belongs to the image in the
first line of \textup{(D.46.10)}.

For clarity, the passage from these local tests to the
induced centralizer follows from the multiplication in
\textup{(D.46.5)}. On a two-cycle of \(\tau\), choose
one coordinate of a lift of a commuting element. Its
partner coordinate is forced by the commutation equation;
consistency of the second equation is precisely commutation
of the chosen coordinate with \(\tau^2\).
The three tests in \textup{(D.46.10)} supply those coordinates.
Their independent corrections are available because the
kernel is the entire \(V^6\). A general element of
\(C_D(\tau)\) has \(\mathfrak S_3\)-projection either
\(1\) or \(\bar\tau\); multiplication by \(\tau^{-1}\)
reduces the latter to the pure case just checked.
We have proved the three actual epimorphisms
\[
 C_E(s)\twoheadrightarrow C_D(s),\quad
 C_E(r)\twoheadrightarrow C_D(r),\quad
 C_E(\tau)\twoheadrightarrow C_D(\tau).
 \tag{D.46.11}
\]
Their kernels are respectively the three spaces in
\textup{(D.46.9)}.

\paragraph{The complete generating fibre.}
Apply the explicit product equation D.23.8 to
\((s,r,\tau^{-1})\). Since \(W_D=0\), its correction
map from \(W^3\) onto \(W\) is surjective for every
prescribed triple of lower conjugators.
Its solution set in conjugators has size \(|W|^2\).
Dividing only by the three kernel centralizers gives
\[
 \frac{|W|^2}{|W^s||W^r||W^\tau|}=|W|.
 \tag{D.46.12}
\]
The epimorphisms \textup{(D.46.11)} show that allowing
all lower centralizer choices gives the same set, rather
than enlarging it: every such change has a lift already
centralizing the corresponding upper peripheral element.
Thus \textup{(D.46.12)} is also the full ordinary
conjugator-union count.

Every solution generates \(E\). Its image surjects onto
\(D\), and its intersection with the irreducible \(W\)
is either \(W\) or zero. In the latter case the image
would be a complement isomorphic to \(D\). It cannot
contain the prescribed conjugate of \(\tau\), whose
upper order \(2\ell^2p\) exceeds its lower order
\(2\ell p\). This excludes the complement case.
The stabilizer of such a generating map under kernel
conjugation is \(W^D=0\). The count in
\textup{(D.46.12)} is therefore exactly one free orbit.
Fixing \(\tau\) cuts that orbit out by \(W^\tau\),
and yields exactly \(\ell^{\ell a}\) maps.

\paragraph{The pure rigidity that proves actual kernel preservation.}
We need an actual quotient automorphism, so we next prove
the source-kernel comparison.
Consider all peripheral epimorphisms \(F\to H\) above
the identity source map \(q_P\). A conjugate of \(X\)
over \(c\) has translation vector whose coordinate sum
is one of the \(p\) values \(\zeta^b\); every vector
with such a sum occurs. A conjugate of \(Y\) over
\(n_0\) has arbitrary translation in \(k e_0\).
Every lift of \(q_P(z)\) is a conjugate of \(Z\),
because \(1-Z\) is invertible. Thus their total number
is \(p|k|^\ell\), and every one generates: its
\(\ell\)-th power at the first cusp gives a nonzero
constant vector, whose \(P\)-orbit spans \(V\).

These maps are exactly one conjugacy orbit under
\(\pi^{-1}Z(P)\). Conjugation by \(V\) first arranges
the first translation to be \(\zeta^b e_0\); the
constant \(N\)-element removes \(\zeta^b\).
The remaining \(V^c=k\sum_i e_i\) conjugation changes
the second translation by \((1-\zeta)k e_0\), and
therefore sets it to zero. The centre of \(H\) is
trivial, so this also accounts for the entire orbit.
After \(X\) is fixed literally, the orbit is exactly
the free \(V^c\)-orbit. In particular every such source
map has precisely the kernel of \(q_H\).

Let now \(\phi:\widehat B\to E\) be any normalized
solution in the preceding fibre, so \(\phi(\tau)=\tau\).
Its conjugator \(f=s\phi(s)\) projects to the given
element of \(D_F'\). Since \(W\subset E_F'\),
it belongs to \(E_F'\). The literal identities of D.44
give \(\phi(x)=x\), \(\phi(y)=f^{-1}yf\), and
positive peripheral conjugacy at \(z\).
The image of \(f\) in \(P'\) belongs to the abelian
subgroup \(N\), which centralizes \(n_0\).
Let \(\operatorname{pr}_0:E_F\twoheadrightarrow H\)
be the identity-coordinate projection of the actual
induced image. Then
\(\xi=\operatorname{pr}_0\phi|_F:F\to H\) lies over the identity
map to \(P\). The pure rigidity just proved gives
\(\xi=\operatorname{Inn}(h)q_H\).

Write \(L=\ker q_E\), the normal core of \(\ker q_H\).
For \(w\in L\) and a coset representative \(t\),
the element \(twt^{-1}\) belongs to \(\ker q_H\).
Moreover \(\phi(t)q_E(t)^{-1}\in E_F\). Therefore
the identity-component projection of
\(q_E(t)\phi(w)q_E(t)^{-1}\) is a conjugate of
\(\xi(twt^{-1})=1\).
All six components of \(\phi(w)\) vanish, proving
\[
                L\subset\ker\phi.
 \tag{D.46.13}
\]
Thus every normalized solution factors through this
same \(E\), and is an automorphism because it is onto.
Conjugating back by \(W\) proves the assertion for the
whole peripheral fibre as well.

\paragraph{The old map and the exact common source.}
The source \(q_D\) kills \(K\), by
\textup{(D.46.6)} in every coordinate, so \(\Gamma\to D\)
is an actual epimorphism. A D.44 normalized automorphism
\(\alpha\) of \(\Gamma\) induces the identity on
the identity-component quotient \(P\): its action is
\(x\mapsto x,y\mapsto f^{-1}yf\), and the image of
\(f\in F'\) lies in \(P'\subset N\), centralizing
\(n_0\). It therefore preserves \(\ker q_P\) and
its six-coset normal core, inducing an actual
\(\alpha_D\in\operatorname{Aut}(D)\).

Residuals commute with continuous epimorphisms, by the
exact proof D.43.7. Thus the projection of \(q_E(K)\)
to the identity component is
\(O^p(O^\ell(H))=V\).
On the other hand \(q_E(K)\subset W\), since
\(D\) is a quotient of \(\Gamma\).
It is a nonzero \(D\)-submodule of the irreducible
\(W\), so
\[
                         q_E(K)=W.
 \tag{D.46.14}
\]
The actual source image in \(\Gamma\times E\)
therefore contains the entire kernel \(1\times W\)
and projects onto \(\Gamma\). This proves
\textup{(D.46.4)} and identifies \(D\) as the full
common quotient for these two source maps.

Each of the \(\ell^{\ell a}\) normalized
automorphisms \(\beta\) of \(E\) over
\(\alpha_D\) gives the automorphism
\[
               (g,e)\longmapsto(\alpha(g),\beta(e))
                   \quad\hbox{of }\Gamma\times_D E.
 \tag{D.46.15}
\]
Its source map generates the whole fibre product.
For prescribed old peripheral conjugators, choose the
upper conjugators with their matching \(D\)-images,
which is permitted by the complete fixed-conjugator
calculation above; their pairs are actual conjugators
in \(J\). Also \([J_F,W]=W\) makes
\(J_F^{\rm ab}\to\Gamma_F^{\rm ab}\) an isomorphism.
Thus \(s\sigma(s)\) lies in \(J_F'\) and projects
to the literal old conjugator. Conversely every
normalized peripheral source map to \(J\) above
the specified \(\alpha\) gives one of the just
classified maps to \(E\); \textup{(D.46.13)}
then proves that it is exactly a map
\textup{(D.46.15)}. This proves the full fibre count,
including source generation and factorization.

The source kernel in \textup{(D.46.4)} contains
\(O^\ell(K)\), because its quotient on \(K\)
is the elementary abelian \(\ell\)-group \(W\).
It is strictly smaller than \(K\), by
\textup{(D.46.14)}. The claimed strict return comparison
is therefore the actual quotient diagram
\[
 \widehat B/O^\ell(K)\twoheadrightarrow J
       \twoheadrightarrow\Gamma,
 \qquad \ker(J\to\Gamma)=W\ne1.
 \tag{D.46.16}
\]
Its pure source kernel is contained in \(O^\ell(F)\),
so the original \(P_\ell\) is a quotient of \(J_F\).
Every finite \(\ell\)-quotient of \(J_F\), composed
with \(F\to J_F\), factors through that same
\(P_\ell\). This proves its asserted maximal
pro-\(\ell\) property even though the new kernel
itself is an \(\ell\)-group.

Finally, for each of the normalized conjugator images
\(f_J\in J_F'\), choose one preimage
\(\widetilde f\in F'\). Its entire word fibre is
\[
       \widetilde f\,
         \ker\bigl(F'\longrightarrow J_F'\bigr).
 \tag{D.46.17}
\]
The epimorphism on closed commutator groups follows from
the actual epimorphism \(F\to J_F\).
This states the word-to-image map separately from the
finite number of image parameters.

For adjacent repetition, \(L/O^p(O^pL)\) is an
extension of the two pro-\(p\) groups
\(O^pL/O^p(O^pL)\) and \(L/O^pL\), hence is
pro-\(p\). The defining maximal quotient and the
reverse subgroup inclusion give
\(O^p(O^pL)=O^pL\).
This identity and \textup{(D.46.16)} specify both
comparisons without asserting a lift to the entire
return cover beyond the proved quotient \(J\).
\end{proof}

The literal model program is
\texttt{bc\_return\_prime\_model.py}.
Its full centralizer companion is
\texttt{bc\_return\_prime\_centralizers.py}.
Together they construct the six-coset maps and retain actual word
preimages for the whole new kernel, solve the complete
peripheral product map, and lift a complete generating
set of each base centralizer. At \((\ell,p)=(5,7)\),
the new dimension is \(180\), the three invariant
dimensions are \(90,60,30\), and the base and upper
\(\tau\)-orders are \(70,350\). The full literal
fibre has \(5^{180}\) maps and its normalized fibre
has \(5^{30}\). These are finite tests with all source
and conjugator choices retained; the uniform proof
above supplies every distinct pair \(\ell,p\geq5\).
